\chapter{Introduction}
\label{ch:introduction}
% The introduction motivates your work and puts it into a bigger context.
% It should answer the following questions:
% What is the background of this work?
% What is the state of the art?
% Why is this project necessary to advance the state of the art?
% What are the problems that have to be solved and why are they difficult?
% What are your contributions to solve these problems?
% How do you evaluate your solution to show that it is adequate and applicable?

% An introduction written along these questions naturally follows the \textit{\gls{spse}}\footnote{%
%   The \acrshort{spse} approach was established in a book~\cite{Hoey83}, but is also briefly summarized in a more recent article~\cite{MP12}, which is available online.
% } approach.
% In the \emph{situation}, you set the scene for your work and catch the interest of the readers by showing the importance and generality of the scene.
% In the \emph{problem}, you spot an issue in the scene and show why and how it significantly taints the scene.
% In the \emph{solution}, you outline your solution to that issue.
% Finally, in the \emph{evaluation}, you present the main arguments why the claimed solution actually does solve the problem.

% In the following chapters, you will elaborate each of the four \gls{spse} elements in detail:
% In \textsl{Background}, you lay the foundations for an in-depth understanding of the situation and the problem.
% In \textsl{Related Work}, you show how others have address this (or similar) problems and why their solutions are not sufficient or applicable.
% In \textsl{Implementation}, you specify your solution, which you then evaluate rigorously for strengths and weaknesses in \textsl{Results}.

% At the end of the introduction, you should explicitly show this structure to the reader by briefly explaining how this report is organized.
% Instead of using the general \gls{spse} terminology and the chapter names mentioned above, we urge you to use the domain-specific terminology established in the introduction and point to chapters using cross references (e.g., refering to \cref{ch:background} instead of ``the Background chapter'').


Over the past decades, music production has moved almost entirely into the digital domain. Recording, editing, mixing, and effect processing are today mostly carried out inside a \acrfull{daw}, where a large ecosystem of software effects replaces what used to require dedicated analog hardware.
% ~\cite{zoelzer2011dafx}
The same is now happening on stage: digital amplifier simulators and multi-effects processors such as the Line~6 Helix~\cite{line6helix}, the Kemper Profiler Stage~\cite{kemperprofiler}, the Neural~DSP Quad Cortex~\cite{neuraldspqc}, and the IK~Multimedia TONEX~\cite{iktonex} process the guitar signal in real time and are increasingly used as replacements for analog amplifiers and pedalboards. Projects such as the Neural Amp Modeler~\cite{atkinson_nam} further allow musicians to create digital captures of their own equipment. This illustrates a demand for methods that identify and faithfully reproduce the behavior of analog guitar gear.\\

Guitar effects cover a wide range of signal processing operations, among them distortion effects (overdrive, distortion, fuzz), dynamics processing (compressor, noise gate), filtering (equalizer, wah), modulation effects (chorus, flanger, phaser, tremolo, vibrato), pitch effects (octaver, harmonizer, pitch shifter), as well as time-based effects such as delay and reverb~\cite{zoelzer2011dafx}.  Nonlinear effects remain challenging, as their characteristic sound depends on the behavior of analog nonlinear circuitry. This work focuses only on the identification and reproduction of \emph{distortion}, which is one of the most defining effects of the electric guitar.\\


The digital emulation of analog audio devices is known as \acrfull{va} modeling. Approaches to \gls{va} modeling are commonly classified as white-box, grey-box, or black-box~\cite{vanhatalo2022review, comunita2025differentiable}. White-box models are derived from the circuit schematic and simulate its behavior using physical models of the individual components~\cite{yeh2009thesis, wright2020realtime}. They can be highly accurate, but require full knowledge of the circuit. Black-box models, on the other hand, require no knowledge of the internal circuitry and are identified purely from input--output measurements of the device~\cite{wright2020realtime}. 
Grey-box models combine both, using a simplified, circuit-informed structure whose parameters are fitted to measured data~\cite{comunita2025differentiable}.\\


Nonlinear black-box modeling is not unique to audio. 
In wireless communications, \acrfullpl{pa} exhibit nonlinear distortion, and behavioral models are often used to describe them for \acrfull{dpd}~\cite{ghannouchi2015behavioral}. 
Such behavioral models are usually identified from input--output observations alone and are therefore black-box models~\cite{silveira2009parallel}.\\

While single Wiener and Hammerstein structures have been applied to guitar distortion~\cite{novak2010synchronized, eichas2016blackbox, eichas2018wiener}, the \gls{pwh} model and its identification from multisine measurements have, to the best of our knowledge, so far only been used for \glspl{pa}. Transferring them to audio is not straightforward: a \gls{pa} is modeled in a narrow band around its carrier, whereas a guitar signal spans about ten octaves, and a distortion pedal clips much harder and depends strongly on the input level. At the same time, the model consists only of a few filters and static nonlinearities, which makes it small enough to run in real time on an embedded platform.\\

In this work, we therefore apply the \gls{pwh} model to the identification
and real-time emulation of guitar distortion effects. The work is divided into
two main tasks:
\begin{enumerate}
  \item \textbf{Identification:} We adapt an existing \gls{pwh} identification
        codebase (\cref{sec:related_pwh}), originally developed for
        \glspl{pa}, to audio signals and use it to identify models of several
        digital and analog guitar distortion effects from input--output
        measurements.
  \item \textbf{Real-time implementation:} We implement the identified
        \gls{pwh} models on the Graetli (a real-time \gls{dsp} guitar pedal
        platform introduced in \cref{ch:background}).
\end{enumerate}

We evaluate the identified models by comparing their outputs with recordings
of the original effects in terms of spectral distance for a range of input
signals. The real-time implementation is assessed by measuring the processing
time per audio block on the Graetli and comparing it to the available
real-time budget. Finally, the original pedal and the Graetli are compared
side by side by processing the same input signals on both and again
evaluating the spectral distance.\\


This report is structured as follows. \Cref{ch:background} introduces the
\gls{pwh} model and the Graetli platform. \Cref{ch:related_work} reviews
related approaches to distortion modeling and explains how they relate to
this work. \Cref{ch:implementation} describes the implementation of the
\gls{pwh} identification for audio, the problems encountered, and the
real-time implementation on the Graetli. The validation and performance
evaluation are presented in \cref{ch:results}. Finally,
\cref{ch:conclusion} summarizes the contributions of this work and discusses
directions for future improvements.